\documentclass[11pt,a4paper]{article}
\usepackage[T1]{fontenc}
\usepackage[utf8]{inputenc}
\usepackage{lmodern,amsmath,amssymb}
\usepackage[margin=25mm]{geometry}
\usepackage[hidelinks]{hyperref}
\hypersetup{pdftitle={The cost of a mark: a Newton-age floor for the Galileo comparison},pdfauthor={navstokgap research working notes}}
\newcommand{\tightlist}{\setlength{\itemsep}{0pt}\setlength{\parskip}{0pt}}
\setlength{\emergencystretch}{3em}
\setcounter{secnumdepth}{0}
\begin{document}
\hypertarget{the-cost-of-a-mark-a-newton-age-floor-for-the-galileo-comparison}{%
\section{The cost of a mark: a Newton-age floor for the Galileo
comparison}\label{the-cost-of-a-mark-a-newton-age-floor-for-the-galileo-comparison}}

Newton takes the sagitta and the swept area of the Galileo comparison to
zero in Lemmas X--XI and keeps only their ratio to the time. That limit
is a statement about the unmarked curve, and it has no floor: every
partition of the fall carries impulses, sagittas and areas that vanish
with the mesh (Proposition 1). A floor appears as soon as the comparison
must be \emph{recorded}, and it appears from three premises that
Newton's own optics supplies: light is made of least corpuscles of
impulse \(p\) (M1), a mark made with light of one colour fixes a
position no better than the interval of fits \(\Lambda\) (M2), and a
mark that fixes a position within \(\delta\) leaves the impulse it
delivered undetermined within \(\Lambda p/\delta\) (M3). Under M1--M3
the falling parabola can be told from the inertial line over a duration
\(\tau\) only if

\[\boxed{\tau\,\Delta E=\frac{F^2\tau^3}{2m}>8\Lambda p,\qquad
A_{\rm inertial,fall}=\frac{vF\tau^3}{6m}>\frac{8v\Lambda p}{3F},}\]

and Democritus insertion of marks into the trajectory stops at the mesh
\(\tau_*=(16m\Lambda p/F^2)^{1/3}\): finer marks record free motion only
(Theorems 1--2). Dropping M3 and keeping Newton's determinate fits gives
an exact countermodel with floor zero (Section 5), so M3 is the single
premise that carries \(h>0\); it is a statement of indeterminacy, and no
Newton-age mechanics or optics derives it. Newton measured
\(\Lambda=1/89000\) inch; he could not measure \(p\); and his corpuscle
sizes make \(\Lambda p\) colour dependent. On the modern identification
\(\Lambda p=h/2\) the floor is \(4h\). This is N02's principle,
obstruction and countermodel. Two later notes remove its two weaknesses.
The \href{planck-gap-derivation.md}{derivation note} removes the
protocol dependence: for every finite protocol of marks with
\(\delta\Delta\ge\kappa\) the comparison needs \(F^2\tau^3>9m\kappa\),
and two marks suffice above \(36m\kappa\). The
\href{mark-cost-and-statistical-floor.md}{mark-cost note} identifies M3
and fixes its constant: for a mark made by coupling the body's position
to a probe, the error operator and the impulse delivered are canonically
conjugate, so \(\kappa\ge\hbar/2\) exactly by Robertson's inequality,
and M3 is that inequality stated in Newton's two quantities. In the
Gaussian statistical model the floor becomes
\(\tau\Delta E\ge48z_{1-\epsilon}^2\kappa\), sharp, for marks whose
error and recoil are uncorrelated; a correlation \(\rho\) lowers the
effective cost to \(\frac\hbar2\sqrt{(1-\rho)/(1+\rho)}\), which is the
one resource that moves the floor. An earlier version of this paragraph
reported \(\kappa\ge2\hbar\) from the derivation note's first
Proposition 6; that value was too large by four and invalidly proved,
and the worst-case widths used here have no non-vacuous quantum instance
at all. Exploratory; no ledger promotion.

\hypertarget{what-newtons-limit-commits-to}{%
\subsection{1. What Newton's limit commits
to}\label{what-newtons-limit-commits-to}}

Take \(m,F,v,\tau>0\), horizontal coordinate \(x\), downward coordinate
\(y\), inertial path \(q_I(t)=(vt,0)\) and falling path
\(q_F(t)=(vt,Ft^2/2m)\) on \([0,\tau]\), as in
\href{newton-insertion-action.md}{N01}. The sagitta, the inertial--fall
area and the gained energy are

\[s(\tau)=\frac{F\tau^2}{2m},\qquad A(\tau)=\frac{vF\tau^3}{6m},\qquad
\Delta E(\tau)=\frac{F^2\tau^2}{2m},\qquad \frac{3F}{v}A=\tau\Delta E .\]

Lemma X says \(s\propto\tau^2\) at the start of the motion; Lemma XI,
Corollaries 4--5, that the curved segment is one third of the tangent
triangle and both scale as \(\tau^3\); Proposition VI reads the force
off the limit \(2ms/\tau^2\); Proposition I builds the orbit as a
polygon with ``a single but great impulse'' at each vertex and lets the
number of vertices be ``augmented in infinitum'' (Motte/Wilkins,
passage; the \href{principia-constant-force-action.md}{Principia note}
has the paragraph map). The impulses, sagittas and areas at the vertices
all vanish in that limit; the force survives as a ratio.

\textbf{Proposition 1 (no geometric floor).} For every partition
\(0=t_0<t_1<\dots<t_n=\tau\) with steps \(\tau_j\) and mesh
\(|\pi|=\max\tau_j\), the vertex impulse \(F\tau_j\), the step sagitta
\(F\tau_j^2/2m\) and the step area \(vF\tau_j^3/6m\) tend to zero with
the mesh, their sum obeys \(\sum_jvF\tau_j^3/6m\le vF\tau|\pi|^2/6m\),
and the ratio \(2m\,s(\tau_j)/\tau_j^2=F\) is exact at every step. No
partition, however fine, produces a contradiction.

\emph{Proof.} Direct from the formulas; the bound on the sum is C027's
computation with the constant-force chord replaced by the inertial
tangent. \(\square\)

Berkeley's objection that evanescent increments are ``neither finite
Quantities nor Quantities infinitely small, nor yet nothing''
(\href{../docs/classics/Berkeley_Analyst_1734_Wilkins.md}{Analyst §35},
passage) is answered inside the geometry by the limit ratio, as
C027--C029 and the full-state classical cuts answer it inside mechanics.
A floor therefore has to come from what it takes to \emph{mark} the
curve, since refinement consistency of the curve offers none. That is
the reading of ``Democritus insertion'' used here: inserting a point
into the trajectory means making a physical mark on the body at a time,
and marks have a cost.

\hypertarget{the-comparison-as-a-record}{%
\subsection{2. The comparison as a
record}\label{the-comparison-as-a-record}}

A \textbf{mark} is an interaction between the body and a recording
system that leaves a datum about the transverse coordinate \(y\) at a
time \(t\). Its \textbf{resolution} \(\delta\) is the width of the
interval to which the datum confines \(y(t)\). Its \textbf{recoil width}
\(\Delta\) is the width of the interval to which the totality of
records, made before or after the mark, confines the transverse impulse
the mark delivered to the body. Both are worst-case widths; nothing
probabilistic is assumed.

\textbf{Two-mark protocol.} A preparation mark at \(t=0\) with
\((\delta_0,\Delta_0)\) and a terminal mark at \(t=\tau\) with
resolution \(\delta_1\). From the preparation the inertial prediction is
an interval \(I_I\) of width

\[w=\delta_0+\frac{\Delta_0\tau}{m},\]

and the falling prediction is \(I_F=I_I+s(\tau)\). The terminal datum
decides between them in the worst case exactly when

\[s(\tau)>w+\delta_1. \tag{1}\]

The velocity \(v\) plays no role in (1); it enters only the conversion
of the action-scaled product \(\tau\Delta E\) into the geometric area
\(A\).

\textbf{Insertion protocol.} Marks at \(t_j=j\tau/n\), each with
\((\delta_j,\Delta_j)\). The force is \emph{recorded on step \(j\)} when
the comparison restricted to \([t_j,t_{j+1}]\), with mark \(j\) as
preparation and mark \(j+1\) as terminal, satisfies (1) with \(\tau\)
replaced by \(\tau_n=\tau/n\).

\hypertarget{three-premises-from-newtons-optics}{%
\subsection{3. Three premises from Newton's
optics}\label{three-premises-from-newtons-optics}}

All four Opticks passages below are stored with their line anchors in
the
\href{../docs/classics/Newton_Opticks_1730_fits_and_queries.md}{source
companion}, which prints them verbatim from the 1730 fourth edition.

\begin{itemize}
\tightlist
\item
  \textbf{(M1) Least impulse.} Light consists of corpuscles (``Are not
  the Rays of Light very small Bodies emitted from shining
  Substances?'', Opticks Query 29, passage). Corpuscles of one colour
  are alike; each carries an impulse \(p>0\), and no part of the light
  of that colour carries less. This is Democritus' least part applied to
  the probe.
\item
  \textbf{(M2) Least length.} A mark made with light of one colour has
  resolution \(\delta\ge\Lambda\), the interval of fits of that colour.
  Newton defines the interval as ``the space it passes between every
  return and the next return'' of the disposition to be reflected (Book
  II Part III, Definition after Prop. XII, passage), and measures it:
  ``the Intervals of their Fits of easy Reflexion are the 1/89000th part
  of an Inch'' for the yellow--orange confine (Prop. XVIII, passage),
  that is \(\Lambda\approx0.285\,\mu\)m. A datum obtained through the
  fits repeats every \(\Lambda\) and so fixes a length no better than
  \(\Lambda\); Newton's rings are exactly such a gauge.
\item
  \textbf{(M3) Mark trade-off.} A mark made with light of one colour
  that fixes the position within \(\delta\) has recoil width
  \(\Delta\ge\Lambda p/\delta\). Newton's theory motivates the
  \emph{pre-mark} half of this: whether a corpuscle is reflected or
  transmitted at the body is settled by its fit, the fit depends on the
  path length modulo \(\Lambda\), and that path length is the very
  position the mark is to find, so no record made before the mark fixes
  the delivered impulse (\(2p\) on reflection, about \(0\) on
  transmission) better than \(p\); confining the ray to a breadth
  \(\delta\) adds the inflexion of Book III (``Do not Bodies act upon
  Light at a distance, and by their action bend its Rays'', Query 1,
  passage; Grimaldi 1665, metadata) with angles of order
  \(\Lambda/\delta\). The premise asserts more: that records made
  \emph{after} the mark do not recover the impulse either. Section 5
  shows this is the whole content.
\end{itemize}

The dimensional gate of
\href{action-unit-dimensional-selection.md}{Q14}, Theorem A, sorts the
premises before any calculation. The marking model's fixed constants are
\(\Lambda\), with dimension vector \((0,1,0)\), and \(p\), with
\((1,1,-1)\). Their unique action product is \(\Lambda p\). A model with
M2 alone (fixed \(\Lambda\)) or M1 alone (fixed \(p\)) holds no action
product and so has floor \(0\) or \(\infty\) for every observable;
Section 5 shows it is \(0\). Both least quantities are needed, and they
are needed as a product, which is what M3 says. The speed of light \(c\)
does not enter the protocol; admitting it would only add the mass
\(p/c\) and a possible factor \(F(mc/p)\) that the argument below never
produces.

\hypertarget{the-floor-and-the-insertion-mesh}{%
\subsection{4. The floor and the insertion
mesh}\label{the-floor-and-the-insertion-mesh}}

\textbf{Theorem 1 (floor for the Galileo comparison).} Under M1--M3,
with the preparation and terminal marks made by light of one colour, the
two-mark protocol decides the comparison in the worst case only if

\[s(\tau)>2\sqrt{\frac{\Lambda p\,\tau}{m}}+\delta_1,\qquad\text{hence}\qquad
\tau\Delta E=\frac{F^2\tau^3}{2m}>8\Lambda p,\qquad
A>\frac{8v\Lambda p}{3F}.\]

\emph{Proof.} By M3, \(\Delta_0\ge\Lambda p/\delta_0\), so
\(w\ge\delta_0+\Lambda p\tau/(m\delta_0)\ge2\sqrt{\Lambda p\tau/m}\),
with equality at \(\delta_0=\sqrt{\Lambda p\tau/m}\). Insert in (1) and
drop \(\delta_1\ge\Lambda>0\): \(F\tau^2/2m>2\sqrt{\Lambda p\tau/m}\).
Squaring, \(F^2\tau^4/4m^2>4\Lambda p\tau/m\), so
\(F^2\tau^3>16m\Lambda p\), which is the displayed bound; the area
follows from \((3F/v)A=\tau\Delta E\). \(\square\)

The constant is protocol dependent, the form is not. A three-mark
protocol with no prepared velocity, marks at \(0,\tau/2,\tau\) and the
mid-chord test \(y(\tau/2)-\tfrac12[y(0)+y(\tau)]\), has deviation
\(s(\tau)/4\); the first recoil cancels in the chord test, the second
enters as \(\Delta\tau/4m\), the three resolutions enter with weights
\(1,\tfrac12,\tfrac12\), and the same minimization gives
\(F^2\tau^3>128m\Lambda p\), that is \(\tau\Delta E>64\Lambda p\). Any
protocol built from marks obeying M1--M3 gives
\(\tau\Delta E>\kappa\Lambda p\) with a positive number \(\kappa\) fixed
by the protocol, because (Theorem A) \(\Lambda p\) is the only action
the model contains. The action-scaled product carries the universal
floor; the geometric area carries the extra factor \(v/F\). This is the
Q14 gate seen from Newton's side, and the reason the anchor's conversion
factor \(3F/v\) is indispensable.

\textbf{Theorem 2 (insertion mesh).} Under M1--M3, in the insertion
protocol the force is recorded on a step only if \(\tau_n>\tau_*\),
where

\[\tau_*=\left(\frac{16\,m\Lambda p}{F^2}\right)^{1/3},\qquad
n<n_*=\frac{\tau}{\tau_*}=\tau\left(\frac{F^2}{16m\Lambda p}\right)^{1/3},\]

and at the finest recording mesh each step's area is
\(A(\tau_*)=8v\Lambda p/3F\). Over \(k\) consecutive steps of a finer
mesh the comparison becomes decidable exactly when \(k\tau_n>\tau_*\).

\emph{Proof.} Theorem 1 applied to one step of duration \(\tau_n\), then
to the duration \(k\tau_n\), using the marks at the ends and ignoring
the interior marks, which the observer may make with \(\delta\to\infty\)
and hence zero recoil. \(\square\)

Below \(\tau_*\) every recorded step is consistent with inertial motion.
The polygon of Proposition I can be refined past \(n_*\), and its
vertices still carry Newton's impulses \(F\tau_n\) in the geometry, but
no mark shows them: the body is free at every recorded instant and
curved over any \(k>n/n_*\) of them. Zeno's arrow and Democritus' cone
(I003, remark 1) receive one answer: adjacent recorded instants are
equal, adjacent recorded slices are equal, and the difference lives in
the unrecorded accumulation over more than \(n/n_*\) of them.

\hypertarget{the-countermodel-determinate-fits}{%
\subsection{5. The countermodel: determinate
fits}\label{the-countermodel-determinate-fits}}

Keep M1 and M2 and take the fits and the inflexion as Newton took them,
as determinate dispositions: the corpuscle's fate at the body is a
function of its phase and its distance from the edge, and the corpuscle
then travels freely to a screen at distance \(D\), where a second mark
of resolution \(\delta_s\ge\Lambda\) records its landing. The
corpuscle's transverse momentum after the mark is then fixed by the two
positions to within \(p(\delta+\delta_s)/D\), and by momentum
conservation so is the impulse it left in the body. Since \(D\) is free,
the recoil width \(\Delta\to0\) at fixed \(\delta\ge\Lambda\): M3 fails.
The two-mark protocol then needs only \(s(\tau)>2\Lambda+\delta_1\),
roughly \(F\tau^2>6m\Lambda\), and

\[\tau\Delta E>\frac{F^2}{2m}\left(\frac{6m\Lambda}{F}\right)^{3/2}
=\sqrt{54\,mF\Lambda^3}\longrightarrow0\quad(F\to0),\]

so the floor over admitted forces and masses is zero, as Theorem A(i)
requires of a model whose only fixed constant is a length. This is the
same closure the classical apparatus results record: position and
momentum records with independent precisions close the canonical product
(C068--C123, \href{action-scale-obstructions.md}{synthesis}), and the
force lens of C070--C071, whose canonical area \(2F^2\ell^3/3m\) is the
same combination as \(\tau\Delta E\), closes with the delay. Theorem 1
is the statement that this lens cannot be resolved below a fixed
multiple of \(\Lambda p\) once M3 holds.

The countermodel isolates the premise. Ignorance of the fit before the
mark, which Newton's theory grants, is removed by the far screen and is
too weak. M3 therefore means that no record, before or after, fixes both
where the corpuscle struck within \(\delta\) and what impulse it left
within \(\Lambda p/\delta\): the delivered impulse is not a function of
the records. In later language this is the complementarity of the
probe's position and momentum records. A Newton-age physicist can state
M3, can see from the fits that it is at least true of all prior records,
and cannot derive it; What Newton printed entails the negation (a
determinate disposition carried by a corpuscle under deterministic
mechanics; he states neither M3 nor its negation). Everything in
Theorems 1--2 follows from M3 with Newton's other materials, and nothing
in them follows without it. This is the precise sense in which the
Newton-age argument for \(h>0\) exists: it is one premise long, and the
premise is the one Newton's determinism excludes.

\hypertarget{newtons-numbers-universality-and-the-quantum-value}{%
\subsection{6. Newton's numbers, universality and the quantum
value}\label{newtons-numbers-universality-and-the-quantum-value}}

Newton had \(\Lambda\) to three figures and had no way to measure \(p\),
so his floor \(\kappa\Lambda p\) is positive and unknown. He also had
the colour dependence of \(\Lambda\): the ring intervals for the utmost
red and utmost violet are ``as 14 to 9'' (Book II Part I, Obs. 13,
passage), and Query 29 makes the red corpuscles the biggest and the
violet the least (passage). If bigger means more impulse, \(\Lambda p\)
increases toward the red and the floor over the spectrum is its violet
value; if size and impulse are unrelated the floor is
\(\min_{\rm colour}\Lambda p\), positive because the spectrum is
bounded. In Q14's terms the fixed constants \(\{\Lambda_i,p_i\}\) give
the action products \(\Lambda_ip_j\) and the dimensionless ratios
\(\Lambda_i/\Lambda_j\), \(p_i/p_j\), so Theorem A(ii) allows any
positive function of them; Newton-age data fix its positivity and not
its value or constancy.

Universality across colours is what Planck's constant adds and what the
Newton-age argument cannot reach. With \(p=h/\lambda\) and
\(\Lambda=\lambda/2\) (the fits alternate every half wavelength),
\(\Lambda p=h/2\) for every colour, and Theorem 1 reads
\(\tau\Delta E>4h\). The quantum benchmark of N01 is the same
combination with a different criterion: perfect discrimination of the
closed-loop phase needs
\(|\Delta S_{\rm loop}|=\tfrac43\tau\Delta E\ge\pi\hbar\), that is
\(\tau\Delta E\ge3h/8\). The factor between \(4h\) and \(3h/8\) is the
difference between worst-case intervals and optimal quantum
discrimination, and between an open two-mark record and a closed
interferometric one. Both say that the inertial--fall lens must carry an
action-scaled product of order \(h\) before it is a distinguishable
alternative.

\hypertarget{where-hto0-and-h0-both-hold}{%
\subsection{\texorpdfstring{7. Where \(h\to0\) and \(h>0\) both
hold}{7. Where h\textbackslash to0 and h\textgreater0 both hold}}\label{where-hto0-and-h0-both-hold}}

Newton's limit lives in the geometry of the unmarked curve, where
Proposition 1 shows it is sound and C027 shows the chord error closes as
\(|\pi|^2\) at every mesh, marks or no marks. The floor lives in the
marks. The two are compatible because the force is a ratio approached
from durations above \(\tau_*\), where the comparison is recordable, and
because Lemma I's ``ultimately equal'' quantities become literally equal
for every record once the mesh passes \(\tau_*\). Below that mesh
Berkeley's departed quantities are the sagitta and the interval of a
step that no mark can distinguish from free motion, and their ratio is
still \(F\) because the geometry does not depend on being marked.

For I003 this places the two remarks. Remark 2, the obstruction to
mapping time to positions as a double limit, is Theorem 1 and needs M3.
Remark 3, an \(h\) that controls the convergence, is Theorem 2 and the
mesh \(\tau_*\): the marked polygon converges to Newton's curve on
meshes above \(\tau_*\) and stops carrying information below it. The
area-rate floor \(GM/2c\) noted earlier in STATE is a rate and gives no
area floor, since \(GM\tau/2c\to0\); it is set aside.

\hypertarget{consequence-for-state}{%
\subsection{8. Consequence for STATE}\label{consequence-for-state}}

N02 asked for one non-quantum consistency principle for joint
time/position refinement of the Galileo comparison, its cut state and
composition, and either a zero-action contradiction or an exact
countermodel. The principle is M3, stated for marks; the cut state at a
mark is the pair of intervals \((\delta,\Delta)\); composition is the
worst-case propagation of Section 2; the contradiction with zero-action
refinement is Theorems 1--2; the countermodel is Section 5. The
principle turns out to be an added indeterminacy premise, statable in
Newton's terms and excluded by what he printed, and it is the entire
difference between a floor and none; Newtonian mechanics and Newton's
optics contain no consistency requirement that yields it. The user's
necessity question is thereby sharpened to a question about M3 alone.

Next: (a) ask whether M3 has a consistency derivation of its own, for
instance from requiring that marks on probes obey the same premise as
marks on bodies without a far-screen loophole, which is the point where
a non-commuting record structure must enter (G07's ``two scale-free
non-commuting structures''); (b) state Theorem 1 for a general force law
and for the C070--C071 lens, where the same \(F^2\tau^3/m\) combination
already appears, so that the floor is a phase-space statement rather
than a parabola statement; (c) done: the
\href{../docs/classics/Newton_Opticks_1730_fits_and_queries.md}{Opticks
companion} now holds Props. XII and XVIII, the Definition of the
interval, Obs. 13 and Queries 1 and 29 verbatim with line anchors.
\end{document}
